\section{The Fix: Changing the Variable Nobody Was Varying}
\label{sec:fix}

\subsection{The unexamined constant}

Every experiment in Sections~\ref{sec:deadlock} and~\ref{sec:shaders} varied
something inside the translation layer or the driver, and every one of them held
the \emph{compatibility distribution} fixed. GE-Proton~10-34 had been chosen at
the start because it worked, and it was never afterwards treated as a variable.

A compatibility distribution bundles its own build of the Direct3D~12 translation
layer. Changing the distribution changes that layer. The question ``does a
different build of the translation layer produce shaders the driver survives?''
had therefore never been asked, despite the entire investigation being about
shaders the driver does not survive.

\subsection{Building a variant cheaply}

Replacing the bundled layer in place would destroy a working configuration. The
approach taken instead was to clone the distribution and replace only the two
libraries of interest.

\begin{lstlisting}
cp -al "$SRC" "$DST"                 # hardlink clone
for d in <four library directories>; do
  rm -f "$DST/$d/d3d12.dll" "$DST/$d/d3d12core.dll"
  cp    "$NEW/d3d12.dll"    "$DST/$d/"
  cp    "$NEW/d3d12core.dll" "$DST/$d/"
done
\end{lstlisting}

A hardlink clone of a 1.4\,GB distribution costs about \textbf{25\,MB}, because
only the replaced files are new data. This is safe \emph{only} because the
replacement is \verb|rm| followed by \verb|cp|---creating a new inode---rather
than an in-place edit, which would modify the original through the shared inode.
That distinction is the entire safety argument and is easy to get wrong.

Two further points, each of which cost a wasted result:

\begin{itemize}
\item \textbf{A prefix upgrade is one-way.} Once an emulated Windows filesystem
      has been upgraded by a newer distribution it cannot be used by an older
      one. Each variant needs its own prefix.
\item \textbf{Seeding a prefix requires more than the prefix directory.} Copying
      only \verb|pfx/| omits the sibling \verb|tracked_files| and \verb|version|
      metadata, and the launcher aborts before the game starts---which reads
      exactly like the variant being broken. This defect made the
      \emph{winning} configuration look like a failure on its first run.
\end{itemize}

\subsection{Result}

\begin{table}[htbp]
\caption{Translation-layer builds, all other variables held fixed}
\label{tab:proton}
\centering
\footnotesize
\begin{tabular}{@{}llrr@{}}
\toprule
\textbf{Build} & \textbf{Layer} & \textbf{Dialogs} & \textbf{GPU} \\
\midrule
Stock (as used throughout) & 2.14.1 & 9 & 0\% \\
Clone with newer layer & \textbf{3.0.1} & \textbf{0} & \textbf{100\%} \\
Clone with same layer & 2.14.1 & --- & exits early \\
Newer distribution & 3.x & 0 & 34\% \\
\bottomrule
\end{tabular}
\end{table}

\textbf{Zero dialogs. The game runs.} The driver's shader compiler does not
crash, because it is no longer being handed the code that crashed it.

The third row is the control: the same clone procedure with the \emph{original}
layer does not fix anything, which rules out the cloning itself as the cause of
the improvement.

\begin{table}[htbp]
\caption{Before and after, discrete GPU}
\label{tab:beforeafter}
\centering
\footnotesize
\begin{tabularx}{\columnwidth}{@{}L{3.0cm}XX@{}}
\toprule
& \textbf{Before} & \textbf{After} \\
\midrule
Assertion dialogs & 9 & \textbf{0} \\
GPU utilisation & 0\% & \textbf{100\%} \\
Video memory & 1{,}960\,MiB, static & 2{,}064--2{,}070\,MiB, live \\
Shader clock & 210\,MHz idle & boost \\
Window & black, indefinitely & title screen at 2880$\times$1800 \\
\bottomrule
\end{tabularx}
\end{table}

\subsection{Verified three independent ways}

Following the project rule that a fix is not a fix until it has been exercised
end to end, the result was confirmed by three methods that do not share a failure
mode:

\begin{enumerate}
\item A \textbf{199-second continuous run} with the kernel driver's own compute
      process list sampled throughout---ground truth from outside the
      compatibility layer, per Section~\ref{sec:graphics1}.
\item A \textbf{screenshot of the in-game graphics menu}, with DLSS listed among
      the upscaling options. That option is only offered when an NVIDIA adapter
      is present, so the game itself corroborates which GPU it is on.
\item A \textbf{screenshot of the title screen} taken from the installed
      launcher command as an ordinary user would invoke it---not from the test
      harness.
\end{enumerate}

The third matters more than it looks. A fix that works only under the harness is
a fix to the harness.

\subsection{The integrated GPU was re-tested, not assumed}

The Radeon path was re-run under the new configuration rather than presumed
intact: \verb|assert=0|, 95\% utilisation. Both GPUs work. The launcher defaults
to the discrete card and accepts a flag to select the integrated one.

\subsection{Launcher design decisions}

Three small decisions in the shipped launcher are worth recording because each
encodes something learned above.

\textbf{Restore the EGL vendor list explicitly.} Per Section~\ref{sec:egl}, the
system-wide pin must be widened to include the NVIDIA vendor, not assumed to be
additive.

\textbf{Detect the panel's native mode from the kernel, not from a
configuration file.} A subtlety: files under \verb|/sys| report a size of zero
regardless of content, so the common shell idiom for ``file exists and is
non-empty'' is always false there. The file must be read and its content checked.

\textbf{Report memory pressure; do not act on it.} The discrete card's memory is
a hard 6{,}141\,MiB budget, and a resident language-model service commonly holds
3{,}576\,MiB of it. The launcher prints a notice when free memory is low. It does
\textbf{not} stop the service, because that service may be doing work the user
cares about more than the game. Tools that silently kill other work to make room
for themselves are a category of behaviour to avoid.

\subsection{Closing the remaining lead}

One lead from the original three-way plan remained: try a newer graphics driver.
It is closed negatively---all available prerelease driver packages resolve to the
same version already installed, 610.57.04. There is no newer driver to try. The
lead is exhausted rather than unexplored, which is a different and more useful
state to record.

\subsection{Undo}

The original distribution is untouched; the variant is a separate directory with
its own prefix. Reverting is deleting two directories and one line in the
launcher. The fix is additive, and the cost of being wrong about it is a few
seconds.
